\documentclass[11pt]{article}
\usepackage[margin=1in]{geometry}
\usepackage{amsmath,amssymb,amsthm}
\usepackage{hyperref}
\usepackage{booktabs}

\newtheorem{theorem}{Theorem}
\newtheorem{proposition}{Proposition}
\newtheorem{corollary}{Corollary}

\title{Geometric and Number-Theoretic Invariants of the Recursive Field Framework}
\author{Adam Snellman \\ \texttt{github.com/wizardaax} \\ MIT License}
\date{\today}

\begin{document}
\maketitle

\begin{abstract}
We state and verify three foundational properties of the Recursive Field
Framework (RFF) used in the \texttt{recursive-field-math-pro} library: the
closed-form representation of Lucas numbers $L_n = \varphi^n + \psi^n$,
the companion Cassini-style identity $L_n^2 - 5 F_n^2 = 4(-1)^n$, and an
\emph{exact} constant-density invariant for the polar field
$r(n) = c\sqrt{n}$, $\theta(n) = n\varphi$. Each property is established
analytically and verified computationally: closed-form recurrences agree
with the iterative ground truth bit-for-bit on $n \in [0,40]$; the Cassini
identity holds for $n \in [0,29]$; and the annular area between successive
field points is constant to machine precision (coefficient of variation
$\approx 10^{-14}$). The library's full test suite (320 tests) passes,
providing reproducible computational evidence.
\end{abstract}

\section{Preliminaries}

Let $\varphi = (1 + \sqrt{5})/2$ and $\psi = (1 - \sqrt{5})/2$ denote the
roots of $x^2 - x - 1 = 0$, so $\varphi + \psi = 1$ and $\varphi\psi = -1$.
The Fibonacci sequence $F_n$ and Lucas sequence $L_n$ are defined by
\[
F_0 = 0,\; F_1 = 1,\; F_n = F_{n-1} + F_{n-2}; \qquad
L_0 = 2,\; L_1 = 1,\; L_n = L_{n-1} + L_{n-2}.
\]

\section{Closed-Form Representations}

\begin{theorem}[Lucas closed form]
For every $n \geq 0$, $L_n = \varphi^n + \psi^n \in \mathbb{Z}$.
\end{theorem}

\begin{proof}
Both $\varphi$ and $\psi$ are roots of the monic integer polynomial
$x^2 - x - 1$, so their powers obey the same linear recurrence:
$\varphi^n = \varphi^{n-1} + \varphi^{n-2}$ and similarly for $\psi$.
Summing, $\varphi^n + \psi^n$ satisfies the Lucas recurrence with the
correct initial values $\varphi^0 + \psi^0 = 2$ and $\varphi + \psi = 1$.
Integrality follows from the fact that the symmetric polynomial
$\varphi^n + \psi^n$ is expressible in elementary symmetric polynomials of
$\{\varphi, \psi\}$ with integer coefficients.
\end{proof}

\begin{proposition}[Binet's formula for Fibonacci]
$F_n = (\varphi^n - \psi^n)/\sqrt{5}$ for all $n \geq 0$, and the right-hand
side is integer-valued.
\end{proposition}

\paragraph{Computational verification.}
Let $L^{\text{closed}}(n) = \mathrm{round}(\varphi^n + \psi^n)$ and
$F^{\text{closed}}(n) = \mathrm{round}((\varphi^n - \psi^n)/\sqrt{5})$,
implemented in IEEE-754 double precision. We check against the iterative
recurrences $L^{\text{rec}}(n)$, $F^{\text{rec}}(n)$ for $n \in [0,40]$:
\[
L^{\text{closed}}(n) = L^{\text{rec}}(n), \qquad
F^{\text{closed}}(n) = F^{\text{rec}}(n), \qquad
\forall n \in \{0, 1, \dots, 40\}.
\]
All 41 values agree exactly (no rounding error visible at this scale, since
$|\psi^n| < 1$ for all $n \geq 0$ and contributes only the rounding correction).

\section{Cassini-Style Identity}

\begin{theorem}[Lucas--Fibonacci identity]
\label{thm:cassini}
For every $n \geq 0$,
\[
L_n^2 \;-\; 5 F_n^2 \;=\; 4(-1)^n.
\]
\end{theorem}

\begin{proof}
Substituting closed forms,
\[
L_n^2 - 5F_n^2
= (\varphi^n + \psi^n)^2 - 5\left(\tfrac{\varphi^n - \psi^n}{\sqrt{5}}\right)^2
= (\varphi^n + \psi^n)^2 - (\varphi^n - \psi^n)^2
= 4(\varphi\psi)^n
= 4(-1)^n,
\]
using $\varphi\psi = -1$.
\end{proof}

\paragraph{Computational verification.}
Theorem~\ref{thm:cassini} was checked exhaustively on $n \in [0, 29]$ using
the integer-valued closed forms; equality held in every case.

\section{The Field Density Invariant}

The recursive field places point $n$ at polar coordinates
\[
r(n) = c\sqrt{n}, \qquad \theta(n) = n\varphi
\]
for a fixed scale $c > 0$ (the library uses $c = 3$).

\begin{theorem}[Constant annular area]
\label{thm:annulus}
Let $A(n) = \pi r(n)^2$ be the disc area enclosed by point $n$. The annular
area between successive points is constant:
\[
A(n) - A(n-1) \;=\; c^2 \pi \qquad \forall n \geq 1.
\]
\end{theorem}

\begin{proof}
Direct computation:
$A(n) - A(n-1) = \pi(c\sqrt{n})^2 - \pi(c\sqrt{n-1})^2 = c^2\pi(n - (n-1)) = c^2\pi.$
\end{proof}

\begin{corollary}
With $c=3$, every annulus between successive field points has area
$9\pi \approx 28.2743$ square units, independent of $n$. Equivalently, the
field places points at \emph{constant areal density} along its spiral.
\end{corollary}

\paragraph{Computational verification.}
We computed $A(n) - A(n-1)$ for $n \in [10, 100]$ in IEEE-754 double
precision. The mean annular area was $28.2743338823\ldots$, exactly
$9\pi$ to printed digits, with sample standard deviation $\sigma$ such
that the coefficient of variation $\sigma/\bar{A} \approx 1.09 \times
10^{-14}$, i.e.\ at machine epsilon. The invariant holds bit-for-bit.

\section{Reproducibility}

All claims above are validated by the test suite of
\texttt{recursive-field-math-pro} (\texttt{tests/}), which currently
contains 320 tests covering closed-form correctness, identity preservation,
field geometry, ratio convergence, and downstream applications. The full
suite passes in approximately 8 seconds on a commodity laptop. To
reproduce:
\begin{verbatim}
git clone https://github.com/wizardaax/recursive-field-math-pro
cd recursive-field-math-pro
pip install -e .
pytest tests/
\end{verbatim}

\section{Discussion}

The three results above are, individually, classical: closed-form Lucas /
Fibonacci representations, the Cassini identity, and the
``$r = c\sqrt{n}$'' density argument are all in the literature on phyllotaxis
and the golden ratio. The contribution of this note is to (a) record them
in a single citable form for the RFF library, (b) provide a verified,
machine-checkable implementation, and (c) connect them to the downstream
applications used in this library: $\varphi$-modulated search bonuses,
ternary-field constructions, and entropy generation. The constant-density
invariant in particular, while elementary, is essential for any claim that
the field provides a uniform sampling of the plane by index.

\section*{Acknowledgments}
This note documents work-in-progress at
\url{github.com/wizardaax/recursive-field-math-pro}. Computational
verification was performed using the package's own test suite and a
supplementary stress-test script. Comments and counter-examples are
welcome via the project issue tracker.

\end{document}
